% !TeX root = ../drafts/05-arithmetization.tex
\section{Arithmetization framework}\label{sec:arith}

The arithmetization is inspired by the M3 model (multi-multiset matching)~\cite{m3}. Rows are unordered, and an instance is accepted exactly when every table's constraints hold and the bus balances.

\subsection{M3 model}\label{sec:m3}

\paragraph{Tables and columns.} A column of height $2^{\kappa}$ is a function $\cube{\kappa}\to\K$. A table is a group of columns of equal height. The arithmetization fixes the number of tables $\ntab$ and, for each table $T_j$, its number of columns $\ncol{j}$, its $\ncons{j}$ constraints and its $\nfl{j}$ bus flushes. The table heights are not fixed: the prover announces each table's log height $\tau_j$ at the beginning of the proof.

\paragraph{Constraints.} The constraints of $T_j$ are $C_{j,1},\dots,C_{j,\ncons{j}}$, each a polynomial of degree at most $d$ ($d=2$ throughout) in the $\ncol{j}$ columns, required to vanish on every row in $\cube{\tau_j}$.

\paragraph{The bus.} All interactions (between tables) share a single \emph{bus}, a channel carrying tuples of up to $m=16$ $\K$-elements (shorter ones are zero-padded). Table $T_j$ has $\nfl{j}$ flushes $\btup_{j,1},\dots,\btup_{j,\nfl{j}}$, each is either a \emph{push} or a \emph{pull}. Each one of the $m$ coordinates in each tuple is a fixed polynomial of degree at most $d$ in the table's columns. A few \emph{boundary} tuples are pushed (resp. pulled) by default to the bus (for instance for the VM state $(\dsep{ST}, \pc, \fp)$, the initial $(\dsep{ST}, \gen^0, \gen^0)$ is pushed and the final $(\dsep{ST}, \gen^{\,\nprog-1}, \gen^{0})$ is pulled).

\paragraph{Balance.} The bus balances when its pushed tuples and its pulled tuples form the same multiset (\ref{def:multiset}).

\paragraph{Domain separation.} The first coordinate of every tuple is a \emph{domain separator}, a fixed constant naming its interaction: $\dsep{ST}=\gen^{0}$ for VM state, $\dsep{MEM}=\gen^{1}$ for memory, $\dsep{BC}=\gen^{2}$ for bytecode. (They are needed since all interractions shares a common bus)

\subsection{Balancing the bus: grand product}\label{sec:gp}

\emph{Fingerprint.} The verifier samples $\alpha=(\alpha_0,\dots,\alpha_3)\in\E^{4}$ and maps each tuple $\btup$ to a single $\E$-element:
\[
\pi_\alpha(\btup)\;=\;\sum_{i=0}^{15}\eq(\alpha,i)\,\btup_i ,
\]
(smaller tuples are padded to $16$ entries with $0$).

\emph{Product check.} Both sides (\emph{pulls} and \emph{pushes}) carry the same number of tuples, and padding each with factors of $1$, which change no product, makes that number a power of two, $2^{\mu}$. The verifier samples $\beta\in\E$, and balance becomes
\begin{equation}
\prod_{\btup\ \text{pushed}}\bigl(\beta-\pi_\alpha(\btup)\bigr)
\;\qeq\;
\prod_{\btup\ \text{pulled}}\bigl(\beta-\pi_\alpha(\btup)\bigr).
\label{eq:gp}
\end{equation}
\emph{Motivation.} For a multiset $P$ of tuples write
\[
  \Pi_P(A,X)\;=\;\prod_{\btup\in P}\bigl(X-\pi_A(\btup)\bigr),
\]
a polynomial in two formal variables with coefficients in $\K$. The two sides of \eqref{eq:gp} are $\Pi_P(\alpha,\beta)$ for $P$ the pushed multiset and the pulled one.

\begin{lemma}[A product identity is a multiset identity]\label{lem:gp}
For finite multisets $P,Q$ of width-$m$ tuples over $\K$, $\Pi_P=\Pi_Q$ if and only if $P=Q$.
\end{lemma}

\begin{proof}
Easy
\end{proof}

\emph{Completeness.} A balanced bus makes \eqref{eq:gp} hold for every $\alpha$ and $\beta$, so completeness is immediate.

\emph{Soundness.} An unbalanced one makes $\Pi_P\neq\Pi_Q$, by the lemma, and \eqref{eq:gp} tests them at the one point $(\alpha,\beta)$, so Corollary~\ref{cor:idtest} gives a soundness error of $5\cdot 2^{\mu}/|\E|$: each of the $2^{\mu}$ factors is of total degree four in $\alpha$ and one in $\beta$, whatever the tuple's width. 

\subsection{Proving the grand product}\label{sec:gkr}

Take one side of \eqref{eq:gp}. Its tuples give its $2^{\mu}$ leaves $v_k=\beta-\pi_\alpha(\btup_k)$, the ones past the last tuple being $1$. Multiply the leaves two by two, then multiply the results two by two, and so on: after $\mu$ layers one value is left, the root $P=\prod_kv_k$. Layer $0$ is the leaves, layer $i$ has $2^{\mu-i}$ nodes, and layer $\mu$ is the root. The prover announces $P$. GKR proves it by walking back down the tree, turning a claim about one layer into a claim about the layer below, until it reaches the leaves and \S\ref{sec:leafstack} takes over.

\subsubsection{Grand-Product argument via GKR}

Write $\widetilde V_i:\E^{\mu-i}\to\E$ for the multilinear extension of layer $i$. Every node of layer $i$ is the product of its two children,
\[
V_i(x)=V_{i-1}(0,x)\,V_{i-1}(1,x),\qquad x\in\cube{\mu-i},
\]
so by Fact~\ref{fact:eq} a claim on $\widetilde V_i$ at a point $r$ is a sum over the hypercube,
\[
\widetilde V_i(r)=\sum_{x\in\cube{\mu-i}}\eq(r,x)\,\widetilde V_{i-1}(0,x)\,\widetilde V_{i-1}(1,x),
\]
which sumcheck (Fact~\ref{fact:sumcheck}) reduces to the two values $\widetilde V_{i-1}(0,\fc)$ and $\widetilde V_{i-1}(1,\fc)$ at one random $\fc$. The prover sends them and a fresh challenge $u$ combines them multilinearly, leaving the single claim $\widetilde V_{i-1}(u,\fc)$: the same kind of claim, one layer lower. The walk starts at the root, where the claim is the announced $P$ itself, and $\mu$ such steps end on the leaves.

\subsubsection{Two layers at a time}

For efficiency, we contract two binary levels at once, a node being the product of its four grandchildren:
\[
V_i(x)=\prod_{a,b\in\{0,1\}}V_{i-2}(a,b,x).
\]
The same sumcheck then reduces $\widetilde V_i(r)$ to the four values $\widetilde V_{i-2}(a,b,\fc)$, which the prover sends and two fresh challenges $(u_0,u_1)$ combine into one claim $\widetilde V_{i-2}(u_0,u_1,\fc)$, two layers down.

If $\mu$ is odd, GKR first processes the root-most binary layer, which has no sumcheck rounds, and then proceeds entirely in radix four.

\subsubsection{Batching several grand-products}

For recursion friendliness, $T$ grand products are proven in one pass, at one sumcheck per layer instead of $T$.

Pad every tree with $1$-leaves up to the depth $\mu$ of the deepest. The verifier then samples a combiner $\lambda \in \E$, and each layer runs a single sumcheck on the random linear combination of the $T$ identities,
\[
\sum_t\lambda^t\widetilde V_i^t(r)=\sum_{x\in\cube{\mu-i}}\eq(r,x)\sum_t\lambda^t\prod_{a,b\in\{0,1\}}\widetilde V_{i-2}^t(a,b,x),\qquad t\in\{0,\dots,T-1\},
\]
drawing a fresh $\lambda$ after every layer, which binds each tree's own tail values and not merely the combination they entered in. A round costs what one tree's round costs; only the layer-end values are per tree. All $T$ trees leave the pass on claims at the one point $\zeta$.

\subsection{Stacking the leaves}\label{sec:leafstack}

The GKR grand-product argument reduces each side of \eqref{eq:gp} to a single evaluation $\widetilde V_0(\zeta)$ of its \emph{leaves}, the product's factors $\beta-\pi_\alpha(\btup)$, one per bus tuple $\btup$.

\paragraph{Flush blocks.} The leaves split into \emph{blocks}, one per flush rule and one per boundary set. A block of a $2^{\kappa}$-row table holds $2^{\kappa}$ leaves; the row-$z$ leaf is $\beta-\sum_{i}\eq(\alpha,i)\,c_{i}(z)$, with $c_{i}(z)$ coordinate $i$ of the row's tuple. Each $c_i$ is a polynomial of degree at most $d$ in the row's columns (\S\ref{sec:m3}).

\paragraph{Layout.} Order the blocks largest first at aligned offsets, padding to $2^{\mu}$ leaves with $1$s, which leave the product unchanged. This is the stacking of \S\ref{sec:stacking} applied to blocks: block $b$, of $2^{\kappa_b}$ leaves, occupies the subcube $\{(z,\mathsf{sel}_b):z\in\cube{\kappa_b}\}$ for a fixed selector $\mathsf{sel}_b$.

\paragraph{Decomposition.} Split $\zeta=(\zeta^{b}_{\mathrm{lo}},\zeta^{b}_{\mathrm{hi}})$ at block $b$'s boundary. Each block is a subcube, so the leaf MLE splits into a public selector times the block's own MLE, plus what the padding contributes:
\begin{equation}\label{eq:gkr_leaves}
  \widetilde V_0(\zeta)\;=\;\sum_b \eq(\mathsf{sel}_b,\zeta^{b}_{\mathrm{hi}})\Bigl(\beta-\sum_{i}\eq(\alpha,i)\,\widetilde c_{b,i}(\zeta^{b}_{\mathrm{lo}})\Bigr)\;+\;\Bigl(1-\sum_b\eq(\mathsf{sel}_b,\zeta^{b}_{\mathrm{hi}})\Bigr).
\end{equation}
That last term is the padding's weight, for the final $\dots 1111$.

\paragraph{Settling it.} The boundary tuples (\S\ref{sec:m3}) and each array's seed and finalization (\S\ref{sec:lookup}) belong to no table. Their entries have degree one: a constant, the index column (\S\ref{sec:idxcol}), a public component, or one committed column. The prover sends, for each such committed multilinear, its evaluation at $\zeta^{b}_{\mathrm{lo}}$. These claims are then deferred to the PCS (\S\ref{sec:stacking}). The verifier finally evaluates each contribution, and subtracts them, with the corresponding padding term, from the claimed value of $\widetilde V_0(\zeta)$. Call the remainder $R_{\mathrm{side}}$: it remains to prove an equation similar to \eqref{eq:gkr_leaves}, summing only on the flush blocks from tables.

Let $\mathcal B$ be the blocks of $T_j$'s flushes in this leaf vector, one per flush rule (\S\ref{sec:m3}). Each holds one leaf per row of $T_j$, so each is $2^{\tau_j}$ tall and all of them split $\zeta$ at the same place, $\zeta^{b}_{\mathrm{lo}}=\zeta_{<\tau_j}$ and $\zeta^{b}_{\mathrm{hi}}=\zeta_{\ge\tau_j}$, differing only through their selector. Exchanging the two sums,
\begin{align*}
  \sum_{b\in\mathcal B}\eq(\mathsf{sel}_b,\zeta^{b}_{\mathrm{hi}})\sum_{x\in\cube{\tau_j}}\eq(\zeta_{<\tau_j},x)\Bigl(\beta-\sum_{i}\alpha^{i}c_{b,i}(x)\Bigr)
  &\;=\;\sum_{x\in\cube{\tau_j}}\eq(\zeta_{<\tau_j},x)\,B_j(x),\\
  B_j&\;=\;\sum_{b\in\mathcal B}\eq(\mathsf{sel}_b,\zeta^{b}_{\mathrm{hi}})\Bigl(\beta-\sum_i\alpha^{i}c_{b,i}\Bigr).
\end{align*}
$B_j$ has degree at most $d$ in $T_j$'s columns, and the verifier builds its coefficients from $\alpha$, $\beta$ and the selectors. What remains to prove is
\[
  \sum_j\sum_{x\in\cube{\tau_j}}\eq(\zeta_{<\tau_j},x)\,B^{\mathrm{side}}_j(x)\;\qeq\;R_{\mathrm{side}} .
\]
Nothing is sent for these blocks. Each constraint of $T_j$ owes a sum of the same shape, so one sumcheck over the tables' rows proves both (\S\ref{sec:air}).

\subsection{Table sumcheck}\label{sec:air}

Number the sides $s=1,\dots,T$, one per grand product (\S\ref{sec:gkr}), and write $B^{s}_j$ and $R_{s}$ for the forms and totals of \S\ref{sec:leafstack}. One sumcheck proves two types of statements:
\begin{align*}
  \sum_{x\in\cube{\tau_j}}\eq(\zeta_{<\tau_j},x)\,C_{j,i}(x)&\;\qeq\;0,\\
  \sum_j\sum_{x\in\cube{\tau_j}}\eq(\zeta_{<\tau_j},x)\,B^{s}_j(x)&\;\qeq\;R_{s},
\end{align*}
the first for every constraint of every table, the second for every side. Both are $\eq$-weighted sums over one table's rows of a polynomial of degree at most $d$ in its columns, which is what lets a single sumcheck carry them all.

\paragraph{Constraints are $\K$-valued.} A constraint $C_{j,i}$ takes values in $\K$, like the columns it reads and like the bus coordinates (\S\ref{sec:m3}); only the $\xi$-powers and $\eq(\zeta,\cdot)$ that weight it lie in $\E$. A relation between machine words is therefore stated lane by lane, three $\K$ constraints rather than one $\E$ constraint, with the tower multiplication unrolled (\S\ref{sec:tab-mul}); the two have the same zero set, since $\{1,y,y^{2}\}$ is a $\K$-basis of $\E$. The prover gains by it: in the round a table joins the batch its columns are still $\K$-valued (below), so a $\K$ constraint is evaluated in $64$-bit arithmetic and meets $\E$ only in the mixed product against its $\xi$-power, while an $\E$ constraint would cost a full $\E$ multiplication per row.

\paragraph{The point is already drawn.} A zerocheck tests its polynomial at a point drawn after the commitment (Definition~\ref{def:zerocheck}), and $\zeta$ is such a point: its coordinates are the challenges of the GKR pass (\S\ref{sec:gkr}), uniform in $\E$ and drawn after the witness was committed. The constraints played no part in producing them. So no fresh point is needed, and the constraints are zerochecked at $\zeta_{<\tau_j}$, where the bus forms already stand.

\paragraph{One challenge.} The verifier samples $\xi\in\E$, after the totals are fixed. Each table mixes everything it owes into one polynomial in its own columns,
\[
  C_j\;=\;\sum_{i=1}^{\ncons{j}}\xi^{\,o_j+i-1}\,C_{j,i}
  \;+\;\sum_{s=1}^{T}\xi^{\,\sigma+s-1}B^{s}_j,
  \qquad
  S_j\;=\;\sum_{x\in\cube{\tau_j}}\eq(\zeta_{<\tau_j},x)\,C_j(x),
\]
with $\sigma=\sum_j\ncons{j}$ the total number of constraints and $o_j=\ncons{1}+\dots+\ncons{j-1}$. The constraint powers are disjoint across tables, so no table's violation is cancelled by another's. The $T$ bus powers are shared by every table, which is what the bus needs, only the totals being pinned. If every constraint vanishes then
\[
  \sum_j S_j\;=\;\sum_{s=1}^{T}\xi^{\,\sigma+s-1}R_{s},
\]
a number the verifier formed before drawing $\xi$. Reaching it forces the $T$ bus equations at once, and a zero in every constraint slot, up to $(n+\sigma+T)/|\E|$ over $\zeta$ and $\xi$ (Corollary~\ref{cor:idtest}), writing $\taumax=\max_j\tau_j$.

\paragraph{One sumcheck.} The $S_j$ are claims of different lengths. Lift each onto $\cube{\taumax}$ by attaching the indicator of the all-ones assignment on the $\taumax-\tau_j$ variables $T_j$ does not use,
\[
  f_j(X_0,\dots,X_{\taumax-1})\;=\;\bigl(C_j\,\eq(\zeta_{<\tau_j},\cdot)\bigr)(X_0,\dots,X_{\tau_j-1})\cdot\prod_{i\ge\tau_j}X_i ,
\]
so that $\sum_{x\in\cube{\taumax}}f_j(x)=S_j$, since $\sum_x\prod_ix_i=1$. With $F=\sum_jf_j$ the whole statement is one claim, $\sum_{x\in\cube{\taumax}}F(x)=\sum_jS_j$, proven in $n$ rounds.

The rounds bind $X_{\taumax-1}$ first and $X_0$ last, so the first $\taumax-\tau_j$ of them touch only variables $T_j$ does not use: there $f_j$ contributes the line $Y\,k\,S_j$, with $k$ the product of the challenges drawn so far. Table $T_j$ enters at round $\taumax-\tau_j$ and binds $X_{\tau_j-1},\dots,X_0$ in turn. A short claim waits, then joins weighted by the rounds it sat out.

The batch ends with one claim per column of $T_j$, at that table's own point $\fc_{<\tau_j}$, each a prefix of the one $\fc$. Those are what the PCS opening settles (\S\ref{sec:stacking}); no column is ever opened at $\zeta$.
